\section{Introduction}
\label{sec:intro}

Prompt engineering is fragile.
A single added space, a swapped separator token, or a rephrased instruction can swing a language model's accuracy by 76 percentage points on the same task \citep{sclar2023quantifying}.
This sensitivity is not eliminated by scaling, few-shot examples, or instruction tuning \citep{mizrahi2024state}, and the performance landscape across prompt formats is highly non-monotonic---format rankings frequently reverse between models \citep{sclar2023quantifying}.
For practitioners deploying language models in production, this fragility poses a serious reliability challenge: without exhaustive evaluation, there is no way to know whether a chosen prompt will fail.

Recent work in mechanistic interpretability has revealed rich geometric structure inside transformer models.
Attention patterns form weighted graphs whose Ricci curvature correlates with model robustness \citep{park2024geometry}.
Attention heads exhibit a U-shaped stability curve across layers, with mid-layer heads being paradoxically the least stable yet most functionally important \citep{bali2026quantifying}.
Internal representations follow constrained belief-update geometry shaped by architectural constraints \citep{piotrowski2025constrained}.
Yet none of this geometric understanding has been connected to the practical problem of predicting which prompts will fail.

On the prompt engineering side, prior work has thoroughly documented the problem \citep{sclar2023quantifying, mizrahi2024state, zhu2023promptbench} but offers no mechanistic explanation for \emph{why} certain formats fail.
Surface-level features like edit distance do not predict degradation magnitude \citep{mizrahi2024state}, strongly motivating the search for deeper structural features.
A predictive geometric signature---computable from a single forward pass---would enable pre-deployment failure screening without expensive evaluation across many inputs.

We bridge this gap by extracting geometric features from attention circuits and testing whether they predict prompt-level performance.
We run \gptsmall on 200 \ssttwo examples across 20 systematically varied prompt formats, extract 23 geometric features from all 144 attention heads, and correlate these features with prompt accuracy.
Our main finding is that six geometric features significantly correlate with prompt performance ($p < 0.05$, permutation-confirmed), with Frobenius norm ($\spearman = +0.474$) and cross-layer variance ($\spearman = -0.472$) as the strongest predictors.
A classifier using these features achieves 60\% leave-one-out accuracy, outperforming both random prediction (50\%) and a logit-confidence baseline that completely fails (0\%).

Our contributions are as follows:
\begin{itemize}[leftmargin=*,itemsep=0pt,topsep=0pt]
    \item We propose using attention circuit geometry---entropy, singular value spectra, Frobenius norms, and cross-layer variance---as predictive features for prompt failure detection, requiring only a single forward pass.
    \item We demonstrate that six geometric features significantly correlate with prompt performance on \ssttwo ($p < 0.05$), with interpretable directions: higher attention magnitude predicts success, while cross-layer inconsistency predicts failure.
    \item We show that a logit-confidence baseline completely fails at distinguishing prompt quality, motivating geometric approaches over output-based heuristics.
\end{itemize}

The remainder of this paper is organized as follows.
\Secref{sec:related} reviews related work on prompt sensitivity, attention geometry, and mechanistic interpretability.
\Secref{sec:method} describes our methodology for extracting and analyzing geometric features.
\Secref{sec:results} presents the experimental results.
\Secref{sec:discussion} discusses implications and limitations, and \secref{sec:conclusion} concludes.
